% One-time imported draft. Editable native source; no Markdown regeneration.
\section{Project contract}\label{005-seven-segment-counter-assignment-count-through-hexadecimal-glyphs-on-digital-lab}

Read a sixteen-byte lookup table and write a requested consecutive range
of its glyphs to the right seven-segment display port. With the default
table, \texttt{first\_digit} = 0 and \texttt{digit\_count} = 16, the
sequence represents hexadecimal indices 0 through F. Each output must be
a one-byte store to address 0xffff0010 in the Default memory layout.

The inputs satisfy 0 \ensuremath{\leq} \texttt{first\_digit} \ensuremath{\leq} 16 and 0 \ensuremath{\leq}
\texttt{digit\_count} \ensuremath{\leq} 16 \ensuremath{-} \texttt{first\_digit}. An empty range,
including \texttt{first\_digit} = 16 and \texttt{digit\_count} = 0,
produces no I/O writes. For other ranges, write the bytes at
glyphs{[}\texttt{first\_digit}{]} through glyphs{[}\texttt{first\_digit}
+ \texttt{digit\_count} \ensuremath{-} 1{]}, once each and in that order. Repeated
table values still require repeated stores. The table bytes may contain
any patterns from 00 through ff; an instructor may replace the default
glyph design.

Perform no writes to another MMIO address, and use byte-sized
transactions so adjacent ports are untouched. Do not write to glyphs,
\texttt{glyph\_guard}, \texttt{first\_digit}, \texttt{digit\_count}, or
\texttt{inputs\_guard}, even temporarily. Produce no console output and
exit normally within 1,000 completed instructions. No host-time delay or
fixed animation rate is required.

\subsection{The device
  interface}\label{005-seven-segment-counter-assignment-the-device-interface}

The source header connects Digital Lab Sim before execution. In this
MARS peripheral model, the right display is at MMIO base + 0x10, the
left display is at +0x11, keypad control is at +0x12, and timer control
is at +0x13. The keypad result is at +0x14. The Default base is
0xffff0000. This project writes only the right display.

A byte is a segment mask, not an ASCII character or a number the device
decodes. The mask table and segment diagram define the intended default
shapes. The left unit is left blank; blank mask 00 is different from the
zero glyph 3f.

\subsection{Guided and free-writing
  versions}\label{005-seven-segment-counter-assignment-guided-and-free-writing-versions}

The guide uses an index and remaining count. An equivalent
implementation can instead compare a current table pointer with an
exclusive end pointer. The grader observes actual execution stores, so
it does not require the reference registers, labels, or branch
arrangement. Transaction address, width, order, count, and values are
part of the published interface and must agree.

The CLI checks the I/O protocol and protected data. It does not draw or
verify the IDE widget. In Debug, stop at the supplied observation
boundaries to inspect individual glyphs. A full-speed Run may reach its
final value before intermediate repaints are seen.

\subsection{Submission}\label{005-seven-segment-counter-assignment-submission}

Submit main.asm and a trace for the first two iterations. Derive the
loop invariant, explain the zero-count path, and distinguish a digit
index from its segment mask. Explain why a final display value alone
cannot establish that a counter visited the required sequence. Run the
public checks and design one custom glyph in a separate copy.
